% eksperymentalne tłumaczenie na włoski

%

\section{Preludium} % lang-pl
\section{Preludio} % lang-it


Panuje nieziemski bałagan, jeśli chodzi o terminologię dla wielościanów lub, ogólniej, wielotopów. By uniknąć ryzyka zgubienia się, podamy jedynie podstawowe definicje, a nie twierdzenia. % lang-pl
Regna una confusione indescrivibile nella terminologia relativa ai poliedri o, più in generale, ai politopi. Per evitare di perderci, presenteremo soltanto le definizioni fondamentali, senza soffermarci sui teoremi. % lang-it


Z każdym wielościanem związana jest grupa jego izometrii: takich przekształceń (odbić, obrotów) bryły na siebie, która zachowuje odległości. % lang-pl
A ogni poliedro è associato il suo gruppo di isometrie: le trasformazioni del poliedro in sé stesso (riflessioni, rotazioni) che preservano le distanze. % lang-it


Tylko odwzorowanie tożsamościowe przenosi flagę na siebie. % lang-pl
Solo l'identità manda una bandiera in sé stessa. % lang-it

\begin{definition}[symetrie wielościanów] % lang-pl
    Wielościan nazywamy izohedralnym (izotoksalnym, izogonalnym), jeśli jego grupa izometrii działa tranzytywnie na zbiorze ścian (krawędzi, wierzchołków). % lang-pl
\begin{definition}[simmetrie dei poliedri] % lang-it
    
    Un poliedro si dice isoedrale (isotossale, isogonale) se il suo gruppo di isometrie agisce transitivamente sull'insieme delle facce (degli spigoli, dei vertici). % lang-it
\end{definition}

\begin{example}
    Wielościany platońskie są izogonalne, izotoksalne i izohedralne jednocześnie. % lang-pl
    
    I poliedri platonici sono contemporaneamente isogonali, isotossali e isoedrali. % lang-it
\end{example}

Edmund Hess, Max Brückner, a później także Branko Grünbaum będą badać wielościany szlachetne, czyli takie, które są izohedralne oraz izogonalne. % lang-pl

Wiadomo, że wielościany platońskie i Keplera-Poinsota są szlachetne; poza nimi także dwusfenoidy i ogólniej sfenoidi. % lang-pl
Edmund Hess, Max Brückner e più tardi anche Branko Grünbaum studieranno i poliedri nobili, cioè quelli che sono sia isoedrali sia isogonali. % lang-it

Si sa che i poliedri platonici e quelli di Keplero-Poinsot sono nobili; inoltre lo sono i disfenoidi e più in generale gli sfenoidi. % lang-it
% https://en.wikipedia.org/wiki/Noble_polyhedron

\begin{example}
    Wielościany Catalana, podwójne ostrosłupy i latawcościany są izohedralne. Bryły dualne do nich, czyli wielościany archimedesowe, graniastosłupy i antygraniastosłupy, są izogonalne. % lang-pl
    
    I poliedri catalani, le bipiramidi e i deltoedri sono isoedrali. I loro duali, cioè i poliedri archimedei, i prismi e gli antiprismi, sono isogonali. % lang-it
\end{example}

Wielościan dualny do izohedralnego jest zawsze izogonalny. % lang-pl
Il duale di un poliedro isoedrale è sempre isogonale. % lang-it


Wielościany foremne mają identyczne foremne ściany i identyczne naroża. % lang-pl
I poliedri regolari hanno facce regolari congruenti e vertici tutti uguali. % lang-it


Wielościany archimedesowe są wypukłe, mają foremne ściany (dwóch lub trzech rodzajów) i są izogonalne, ale nie izohedralne. % lang-pl
I poliedri archimedei sono convessi, hanno facce regolari (di due o tre tipi) e sono isogonali, ma non isoedrali. % lang-it


Wielościany Catalana są dualne do archimedesowych. % lang-pl
I poliedri catalani sono duali dei poliedri archimedei. % lang-it


Wielościany jednorodne mają foremne ściany oraz są wierzchołkowo przechodnie. % lang-pl
I poliedri uniformi hanno facce regolari e sono transitivi sui vertici. % lang-it


Wielościan Johnsona(-Zalgallera) to wypukły wielościan, którego ściany są wielokątami foremnymi. % lang-pl
Un poliedro di Johnson (o di Johnson-Zalgaller) è un poliedro convesso le cui facce sono poligoni regolari. % lang-it


Przez analogię do wielokątów wypukłych mamy: % lang-pl
Per analogia con i poligoni convessi, abbiamo: % lang-it
\begin{definition}[poliedro convesso] % lang-it
    Una regione limitata dello spazio ottenuta come intersezione di semispazi. % lang-it
    La parte ritagliata su una delle superfici piane dalle altre è un poligono, che chiamiamo faccia. % lang-it
    Il lato comune di due facce è uno spigolo, mentre l'estremo comune di due spigoli è un vertice. % lang-it
\end{definition} % lang-it


\begin{definition}[wielościan wypukły] % lang-pl
    Obszar skończony przestrzeni będący przekrojem półprzestrzeni. % lang-pl
I più noti sono le piramidi e i prismi. % lang-it
Oltre a essi parleremo anche di bipiramidi, deltoedri e di vari solidi più o meno regolari. % lang-it
Non daremo invece una definizione di poliedro (non convesso), poiché non esiste una definizione universalmente accettata. % lang-it
\begin{definition}[bandiera] % lang-it
    
    Część wycięta z płaszczyzny przez pozostałe jest wielokątem, nazywamy go ścianą. % lang-pl
    Wspólny bok dwóch ścian to krawędź, wspólny koniec dwóch krawędzi to wierzchołek. % lang-pl
\end{definition} % lang-pl
Najbardziej znanymi są ostrosłupy i graniastosłupy. % lang-pl
Oprócz nich opowiemy także o podwójnych ostrosłupach, latawcościanach oraz różnych mniej lub bardziej foremnych bryłach. % lang-pl
Natomiast nie podamy definicji wielościanu (niewypukłego), ponieważ ma powszechnie akceptowanej tegoż. % lang-pl
\begin{definition}[flaga] % lang-pl
    Zbiór złożony ze ściany, krawędzi oraz wierzchołka takich, że każde leży na poprzednim, nazywamy flagą. % lang-pl
    Una terna costituita da una faccia, uno spigolo e un vertice, tali che ciascun elemento appartenga al precedente, si chiama bandiera. % lang-it
\end{definition} % https://en.wikipedia.org/wiki/Flag_(geometry)

%